\documentclass[conference]{IEEEtran}
\usepackage{cite}
\usepackage{amsmath,amssymb,amsfonts}
\usepackage{algorithmic}
\usepackage{graphicx}
\usepackage{textcomp}
\usepackage{xcolor}
\usepackage{hyperref}
\usepackage{booktabs}
\usepackage{multirow}
\usepackage{array}
\usepackage{float}
\usepackage{listings}
\usepackage{subcaption}
\usepackage{enumitem}
\usepackage{bm}

\lstset{
  basicstyle=\ttfamily\scriptsize,
  breaklines=true,
  frame=single,
  columns=fullflexible,
  captionpos=b
}

\def\BibTeX{{\rm B\kern-.05em{\sc i\kern-.025em b}\kern-.08em
    T\kern-.1667em\lower.7ex\hbox{E}\kern-.125emX}}

\begin{document}

\title{Cocktail Party Speech Separation Using Deep Learning: A Comparative Study of DPRNN-TasNet, Conv-TasNet, and RCNN Architectures}

\author{
\IEEEauthorblockN{Department of Computer Science and Engineering}
\IEEEauthorblockA{Pattern Recognition and Machine Learning (PRML) Course Project}
}

\maketitle

\begin{abstract}
The cocktail party problem, isolating individual speech signals from a single-channel mixture of overlapping speakers, remains one of the most challenging tasks in audio signal processing. Traditional methods such as Independent Component Analysis (ICA) require at least as many microphones as speakers, rendering them inapplicable to the single-microphone setting. This paper presents a comprehensive comparative study of four deep learning architectures for blind source separation: (1)~Dual-Path Recurrent Neural Network with Time-domain Audio Separation Network (DPRNN-TasNet), (2)~Convolutional Time-domain Audio Separation Network (Conv-TasNet) implemented in Python, (3)~Conv-TasNet reimplemented in C++ using LibTorch for production deployment, and (4)~Recurrent Convolutional Neural Network (RCNN) operating on spectral representations. All models are trained on the LibriSpeech corpus with on-the-fly mixture generation and evaluated using scale-invariant signal-to-noise ratio (SI-SNR), signal-to-distortion ratio (SDR), perceptual evaluation of speech quality (PESQ), and short-time objective intelligibility (STOI). We achieve a best SI-SNR of 15.28~dB for 3-speaker separation with the DPRNN-TasNet architecture. The models are integrated into a unified web application featuring real-time separation, spectrogram visualization, and multi-model comparison. We analyze the trade-offs between architectural complexity, computational cost, training stability, and separation quality across speaker counts ranging from 2 to 5.
\end{abstract}

\begin{IEEEkeywords}
Cocktail party problem, speech separation, deep learning, DPRNN-TasNet, Conv-TasNet, RCNN, blind source separation, SI-SNR, permutation invariant training, LibriSpeech
\end{IEEEkeywords}

% ============================================================
\section{Introduction}
\label{sec:intro}

The ability to selectively attend to a single speaker in a noisy, multi-talker environment, colloquially known as the ``cocktail party problem,'' was first articulated by Cherry in 1953. While the human auditory system performs this task effortlessly, computational solutions remain exceptionally difficult, particularly in the single-microphone setting where spatial cues are entirely absent.

Formally, given a single-channel recording $x(t) = \sum_{i=1}^{C} s_i(t) + n(t)$, where $s_i(t)$ denotes the $i$-th speaker's signal and $n(t)$ represents ambient noise, the objective is to recover estimates $\hat{s}_1(t), \hat{s}_2(t), \ldots, \hat{s}_C(t)$ that closely approximate the original source signals. This problem is ill-posed when $C > 1$ since there are infinitely many decompositions of a single observation into $C$ components.

\subsection{Motivation and Applications}

Effective speech separation has broad applications spanning multiple critical domains:

\begin{itemize}[leftmargin=*,nosep]
    \item \textbf{Automatic Speech Recognition (ASR):} Preprocessing noisy multi-speaker recordings to enable accurate transcription in meeting and conference scenarios.
    \item \textbf{Hearing Aids:} Enhancing the target speaker's voice while suppressing interfering talkers for hearing-impaired individuals.
    \item \textbf{Forensics and Law Enforcement:} Isolating individual voices from wiretap recordings, crowd audio, or surveillance feeds for identification and evidence analysis.
    \item \textbf{Teleconferencing:} Separating speakers on overlapping channels to improve call quality and enable speaker-attributed transcripts.
    \item \textbf{Copyright and Content Analysis:} Detecting copyrighted speech or music embedded within mixed audio streams.
\end{itemize}

\subsection{Limitations of Classical Approaches}

Traditional signal processing methods for source separation include Independent Component Analysis (ICA), Non-negative Matrix Factorization (NMF), and Computational Auditory Scene Analysis (CASA). Among these, ICA has been the most widely studied. ICA assumes that sources are statistically independent and seeks a linear unmixing matrix $\mathbf{W}$ such that $\hat{\mathbf{s}} = \mathbf{W}\mathbf{x}$. However, ICA suffers from a fundamental limitation: it requires the number of observations (microphones) to be at least equal to the number of sources, i.e., $M \geq C$. In the single-channel case $M=1$, the system is heavily underdetermined, and ICA fails entirely.

FastICA, a popular variant, achieves separation in multi-channel setups by maximizing non-Gaussianity of estimated sources. However, when applied to single-channel audio, it cannot separate more than one source. Furthermore, ICA makes the assumption of instantaneous mixing, which is violated in reverberant environments where convolutive mixing occurs.

\subsection{Deep Learning Paradigm}

Recent advances in deep learning have fundamentally changed the landscape of speech separation. Neural networks can learn complex, nonlinear mappings from mixed to separated signals by leveraging large corpora of clean speech data. The key insight is that deep models implicitly learn statistical priors over the space of speech signals, enabling them to ``hallucinate'' plausible decompositions even from a single microphone recording.

This paper presents a systematic comparison of four deep learning architectures spanning different design philosophies: time-domain masking (DPRNN-TasNet and Conv-TasNet), spectral masking (RCNN), and production-optimized inference (Conv-TasNet in C++). All models share a common encoder, separator, and decoder paradigm but differ substantially in their separator designs.

\subsection{Contributions}

The principal contributions of this work are:
\begin{enumerate}[leftmargin=*,nosep]
    \item A detailed comparative study of four architectures (DPRNN-TasNet, Conv-TasNet, Conv-TasNet-C++, RCNN) for single-channel speech separation on the LibriSpeech corpus.
    \item Implementation of curriculum learning for multi-speaker generalization (2$\rightarrow$3$\rightarrow$4$\rightarrow$5 speakers).
    \item A production-grade C++ implementation of Conv-TasNet using LibTorch for low-latency inference.
    \item A unified web application integrating all models with real-time spectrogram visualization and multi-model comparison.
    \item Analysis of domain mismatch between synthetic training mixtures and real-world recordings.
\end{enumerate}

% ============================================================
\section{Related Work}
\label{sec:related}

\subsection{Independent Component Analysis}

ICA is a classical statistical method that separates a multivariate signal into additive, statistically independent components. Given observations $\mathbf{x} = \mathbf{A}\mathbf{s}$, ICA recovers $\mathbf{s}$ by estimating the unmixing matrix $\mathbf{W} = \mathbf{A}^{-1}$. The FastICA algorithm achieves this by maximizing the negentropy of estimated sources through fixed-point iteration. While effective for multi-channel signals, ICA requires $M \geq C$ (number of microphones $\geq$ number of sources) and assumes linear instantaneous mixing, both assumptions violated in the single-channel cocktail party problem.

\subsection{Deep Clustering and Permutation Invariant Training}

Hershey et al.\ introduced deep clustering (DPCL), where a neural network maps each time-frequency bin to an embedding space such that bins dominated by the same speaker are close together. Yu et al.\ proposed Permutation Invariant Training (PIT), which resolves the label ambiguity problem by evaluating all possible permutations of source assignments and selecting the one with minimal loss:
\begin{equation}
\mathcal{L}_{\text{PIT}} = \min_{\pi \in \mathcal{P}} \frac{1}{C} \sum_{i=1}^{C} \ell(\hat{s}_i, s_{\pi(i)})
\label{eq:pit}
\end{equation}
where $\mathcal{P}$ is the set of all permutations of $\{1, \ldots, C\}$ and $\ell$ is the per-source loss function.

\subsection{TasNet and Conv-TasNet}

Luo and Mesgarani introduced TasNet, the first fully time-domain separation network, which bypasses the STFT by learning an adaptive encoder-decoder pair via 1D convolutions. Conv-TasNet replaced the LSTM separator with a Temporal Convolutional Network (TCN) consisting of stacked dilated depth-wise separable convolutions, achieving 15.3~dB SI-SNRi on the WSJ0-2mix benchmark with significantly fewer parameters and faster inference than LSTM-based approaches.

\subsection{Dual-Path RNN}

Luo et al.\ proposed the Dual-Path RNN (DPRNN), which addresses the computational bottleneck of processing long sequences by splitting the encoded representation into overlapping chunks and alternating between intra-chunk (local) and inter-chunk (global) recurrent processing. This dual-path strategy reduces the computational complexity from $O(T^2)$ to $O(T \cdot K)$ while maintaining global context modeling, achieving state-of-the-art results of 18.8~dB SI-SNRi on WSJ0-2mix.

\subsection{Spectrogram-Based Methods}

An alternative paradigm operates on Short-Time Fourier Transform (STFT) magnitude spectrograms, predicting time-frequency masks that are multiplied with the mixture spectrogram and inverted via inverse STFT (iSTFT). Methods such as the recurrent convolutional network use an encoder-decoder architecture with a BiLSTM bottleneck to capture temporal dependencies in the spectral domain.

% ============================================================
\section{Dataset Description}
\label{sec:dataset}

\subsection{LibriSpeech Corpus}

All models in this study are trained on the LibriSpeech corpus, a large-scale dataset of read English speech derived from the LibriVox project, an initiative of volunteer-narrated public domain audiobooks. The corpus was originally created by Vassil Panayotov et al.\ at Johns Hopkins University and has since become one of the most widely adopted benchmarks in both automatic speech recognition and speech separation research.

The full LibriSpeech corpus comprises approximately 1,000 hours of read English speech recorded at a 16~kHz sampling rate and encoded as 16-bit FLAC files. The audio is segmented by speaker and organized into subsets based on a quality-of-alignment score between the audio and the corresponding text transcriptions. Speakers cover a broad demographic distribution across gender, age, and accent, providing a rich diversity of vocal characteristics that is critical for training speaker-independent separation models.

We utilize the following subsets in our experiments:
\begin{itemize}[leftmargin=*,nosep]
    \item \textbf{train-clean-100}: Contains approximately 100 hours of clean, high-quality speech from 251 distinct speakers (125 female, 126 male). Each speaker contributes between 15 and 25 minutes of audio on average. This subset serves as the primary training set for all four models due to its manageable size and high transcription alignment quality.
    \item \textbf{train-clean-360}: Contains approximately 360 hours of clean speech from 921 distinct speakers. This larger subset is used by the Conv-TasNet (Python) model for expanded training experiments to study the effect of additional training data on separation performance.
    \item \textbf{dev-clean}: Contains 5.4 hours of speech from 40 speakers who do not appear in the training sets. This subset is used for validation, hyperparameter selection, and early stopping decisions throughout all experiments.
    \item \textbf{test-clean}: Contains 5.4 hours from another disjoint set of 40 speakers, reserved for final performance evaluation and held-out metric reporting.
\end{itemize}

All audio files are single-channel (mono) recordings. During data loading, any files sampled at rates other than the target rate (16~kHz for DPRNN-TasNet, 8~kHz for Conv-TasNet and RCNN) are resampled using high-quality sinc interpolation to ensure consistent input dimensions across the pipeline.

\subsection{On-the-Fly Mixture Generation}

Rather than pre-generating fixed mixtures on disk, our pipeline generates training mixtures dynamically during training. This approach offers several advantages: (1)~effectively infinite data diversity as speaker/segment combinations are sampled randomly each epoch, (2)~no disk storage overhead, and (3)~natural data augmentation through the mixing process itself.

The mixture generation procedure is as follows:

\begin{enumerate}[leftmargin=*,nosep]
    \item \textbf{Speaker Selection:} Randomly sample $C$ distinct speakers from the training set, where $C \in \{2, 3, 4, 5\}$ is the desired number of sources.
    \item \textbf{Segment Extraction:} For each selected speaker, randomly choose a contiguous segment of duration $T_{\text{clip}} = 4$~seconds (64,000 samples at 16~kHz or 32,000 at 8~kHz).
    \item \textbf{SNR Randomization:} Apply a random relative SNR offset $\alpha \in [-5, 5]$~dB to secondary speakers to simulate varying loudness:
    \begin{equation}
    s_i'(t) = s_i(t) \cdot 10^{\alpha_i / 20}, \quad i = 2, \ldots, C
    \end{equation}
    \item \textbf{Additive Mixing:} Sum all sources: $x(t) = \sum_{i=1}^{C} s_i'(t)$.
    \item \textbf{Peak Normalization:} Normalize the mixture to prevent clipping: $x(t) \leftarrow 0.9 \cdot x(t) / \max|x(t)|$.
\end{enumerate}

\subsection{Noise Augmentation with MUSAN}

For noise-robust training, we incorporate the MUSAN (Music, Speech, and Noise) corpus, a freely available dataset curated by David Snyder et al.\ at Johns Hopkins University. MUSAN contains approximately 930 noise clips totaling roughly 6 hours of audio, spanning three categories: (1)~ambient and technical noise recordings such as HVAC systems, street sounds, and office environments, (2)~music segments covering diverse genres, and (3)~babble noise generated from multi-speaker overlaps. All clips are provided as single-channel 16~kHz WAV files.

Noise augmentation is applied to the mixture signal only, while source targets remain clean, at random SNR ratios between 5 and 20~dB:
\begin{equation}
x_{\text{noisy}}(t) = x(t) + \beta \cdot n(t), \quad \beta = 10^{-\text{SNR}/20}
\end{equation}
This ensures the model learns to separate speech from noise-contaminated mixtures while maintaining clean targets for loss computation.

\subsection{Data Splits and Training Configuration}

A ``virtual epoch'' paradigm is employed with a configurable \texttt{epoch\_length} parameter controlling the number of mixtures per training epoch (typically 20,000 to 30,000 for training and 2,000 to 5,000 for validation). Test sets comprise 1,000 to 2,000 held-out mixtures with non-overlapping speakers.

Table~\ref{tab:dataset_config} summarizes the dataset configuration across models.

\begin{table}[t]
\centering
\caption{Dataset Configuration Across Models}
\label{tab:dataset_config}
\begin{tabular}{@{}lcccc@{}}
\toprule
\textbf{Parameter} & \textbf{DPRNN} & \textbf{Conv-TasNet} & \textbf{RCNN} & \textbf{C++} \\
\midrule
Sample Rate & 16 kHz & 8 kHz & 8 kHz & 8 kHz \\
Clip Length & 4 s & 4 s & 4 s & 4 s \\
Train Subset & clean-100 & clean-100/360 & clean-100 & clean-100 \\
Speakers & 2--5 & 2--3 & 2 & 2 \\
Noise Aug. & No & MUSAN & No & No \\
\bottomrule
\end{tabular}
\end{table}

% ============================================================
\section{Methodology}
\label{sec:methodology}

All four architectures share the encoder, separator, and decoder paradigm illustrated in Fig.~\ref{fig:overview}. The encoder transforms the raw waveform (or spectrogram) into a high-dimensional latent representation, the separator estimates per-source masks in the latent space, and the decoder reconstructs the separated waveforms.

\begin{figure}[t]
\centering
\fbox{\parbox{0.95\columnwidth}{\centering
\small
\texttt{Mixture Waveform} \\
$\downarrow$ \\
\texttt{Encoder (1D Conv / STFT)} \\
$\downarrow$ \\
\texttt{Separator (DPRNN / TCN / BiLSTM)} \\
$\downarrow$ \\
\texttt{Mask Estimation ($C$ masks)} \\
$\downarrow$ \\
\texttt{Masked Features = Mask $\odot$ Encoded} \\
$\downarrow$ \\
\texttt{Decoder (1D TransConv / iSTFT)} \\
$\downarrow$ \\
\texttt{$C$ Separated Waveforms}
}}
\caption{Common encoder, separator, and decoder pipeline shared by all models. The separator module and operating domain (time vs.\ spectral) vary across architectures.}
\label{fig:overview}
\end{figure}

\subsection{Loss Function: Permutation Invariant Training with SI-SNR}

All models employ Permutation Invariant Training (PIT) to resolve the label permutation ambiguity inherent in source separation. The per-pair loss is the negative Scale-Invariant Signal-to-Noise Ratio (SI-SNR):

\begin{equation}
\text{SI-SNR}(\hat{s}, s) = 10 \log_{10} \frac{\|\mathbf{s}_{\text{target}}\|^2}{\|\mathbf{e}_{\text{noise}}\|^2}
\label{eq:sisnr}
\end{equation}

where the target projection and noise are defined as:
\begin{align}
\mathbf{s}_{\text{target}} &= \frac{\langle \hat{s}, s \rangle}{\|s\|^2} s \\
\mathbf{e}_{\text{noise}} &= \hat{s} - \mathbf{s}_{\text{target}}
\end{align}

For $C$ sources, the PIT loss evaluates all $C!$ permutations (or uses the Hungarian algorithm for $C > 3$) and selects the optimal assignment:
\begin{equation}
\mathcal{L} = -\max_{\pi \in \mathcal{P}} \frac{1}{C} \sum_{i=1}^{C} \text{SI-SNR}(\hat{s}_i, s_{\pi(i)})
\end{equation}

An optional SNR regularization term with weight $\lambda = 0.1$ is added to encourage general noise suppression.

% ============================================================
\section{Model Architectures}
\label{sec:models}

\subsection{Architecture I: DPRNN-TasNet}
\label{sec:dprnn}

The Dual-Path RNN TasNet architecture addresses the fundamental challenge of modeling both local and global temporal dependencies in long audio sequences. It consists of three stages: a learned encoder, a dual-path recurrent separator, and a learned decoder.

\subsubsection{Encoder}

The encoder replaces the traditional STFT with a learnable 1D convolutional filterbank:
\begin{equation}
\mathbf{w} = \text{ReLU}(\text{Conv1D}(x; N, L, L/2))
\end{equation}
where $N = 64$ is the number of encoder filters, $L = 2$ is the kernel size, and $L/2 = 1$ is the stride. The encoder transforms the raw waveform $x \in \mathbb{R}^{1 \times T}$ into an encoded representation $\mathbf{w} \in \mathbb{R}^{N \times T'}$ where $T' = T/\text{stride}$.

\subsubsection{DPRNN Separator}

The separator first applies group normalization followed by a bottleneck 1$\times$1 convolution projecting from $N$ to bottleneck dimension $B$. The encoded sequence is then segmented into overlapping chunks of size $K = 100$ with 50\% overlap (hop $= K/2$):

\begin{equation}
\mathbf{w}_{\text{chunk}} \in \mathbb{R}^{B \times K \times S}
\end{equation}

where $S$ is the number of chunks. The DPRNN then applies $R = 6$ dual-path blocks, each containing:

\textbf{Intra-chunk RNN (Local Processing):} A bidirectional LSTM processes each chunk independently along the $K$ dimension:
\begin{equation}
\mathbf{h}_k^{\text{intra}} = \text{BiLSTM}_{\text{intra}}(\mathbf{w}_{1:K}^{(s)}), \quad \forall s \in \{1, \ldots, S\}
\end{equation}

The BiLSTM has 2 layers with hidden size equal to the channel dimension and 10\% dropout. A linear projection maps the concatenated forward-backward hidden states back to $B$ dimensions, followed by layer normalization and a residual connection.

\textbf{Inter-chunk RNN (Global Processing):} A second BiLSTM processes across chunks at each frequency position:
\begin{equation}
\mathbf{h}_s^{\text{inter}} = \text{BiLSTM}_{\text{inter}}(\mathbf{w}_{k}^{(1:S)}), \quad \forall k \in \{1, \ldots, K\}
\end{equation}

This enables global context modeling without the $O(T^2)$ complexity of processing the full sequence. Each dual-path block has separate learnable parameters and its own layer normalization with learnable affine parameters $\gamma, \beta$.

\subsubsection{Mask Generation and Overlap-Add}

After the DPRNN blocks, a PReLU activation followed by a 1$\times$1 convolution generates $C$ masks of shape $(B, N \cdot C, K, S)$. Sigmoid activation constrains masks to $[0, 1]$. The overlap-add procedure reconstructs the full-length masked representation:
\begin{equation}
\hat{\mathbf{w}}_c[t] = \frac{\sum_{s} \mathbf{m}_c^{(s)}[t - s \cdot K/2] \cdot \mathbf{w}^{(s)}[t - s \cdot K/2]}{\sum_{s} \mathbf{1}[t \in \text{chunk}_s]}
\end{equation}

\subsubsection{Decoder}

A transposed 1D convolution reconstructs the time-domain signal:
\begin{equation}
\hat{s}_c(t) = \text{ConvTranspose1D}(\hat{\mathbf{w}}_c; 1, L, L/2)
\end{equation}

\subsubsection{Curriculum Learning Strategy}

A multi-phase curriculum trains the model progressively on increasing speaker counts:
\begin{itemize}[leftmargin=*,nosep]
    \item \textbf{Phase 1:} 2-speaker separation (100 epochs, $\text{lr} = 10^{-3}$)
    \item \textbf{Phase 2:} 3-speaker separation (500 epochs, initialized from Phase~1)
    \item \textbf{Phase 3:} 4-speaker separation (500 epochs, initialized from Phase~2)
    \item \textbf{Phase 4:} 5-speaker separation (500 epochs, initialized from Phase~3)
\end{itemize}

Each phase adjusts the output layer to produce $C$ masks and fine-tunes the full network. This strategy helps the model transfer local separation skills learned from simpler tasks to more complex scenarios.

% --- Conv-TasNet ---
\subsection{Architecture II: Conv-TasNet (Python)}
\label{sec:convtasnet}

Conv-TasNet replaces the recurrent separator of TasNet with a fully convolutional Temporal Convolutional Network (TCN), achieving comparable or superior performance with significantly reduced computational cost.

\subsubsection{Encoder-Decoder}

Identical in concept to the DPRNN encoder, Conv-TasNet uses:
\begin{align}
\text{Encoder:} &\quad \text{Conv1D}(1, N{=}512, L{=}16, \text{stride}{=}8) \\
\text{Decoder:} &\quad \text{ConvTranspose1D}(N, 1, L{=}16, \text{stride}{=}8)
\end{align}

The larger filter count ($N = 512$ vs.\ 64) and kernel size ($L = 16$ vs.\ 2) provide a richer learned representation.

\subsubsection{TCN Separator}

The TCN consists of $R = 3$ repeats of $X = 8$ convolutional blocks. Each block contains a depth-wise separable convolution with exponentially increasing dilation:

\begin{enumerate}[leftmargin=*,nosep]
    \item $1 \times 1$ convolution: $B \rightarrow H$ (bottleneck $B{=}128$ to hidden $H{=}512$)
    \item PReLU activation
    \item Layer normalization
    \item Depth-wise convolution: $H \rightarrow H$, kernel $P{=}3$, dilation $2^x$
    \item PReLU activation
    \item Layer normalization
    \item $1 \times 1$ convolution: $H \rightarrow B$ (projection back)
    \item Residual connection from input
    \item Skip connection accumulated to output
\end{enumerate}

The exponential dilation pattern $\{1, 2, 4, 8, 16, 32, 64, 128\}$ yields a receptive field of:
\begin{equation}
\text{RF} = R \cdot P \cdot \sum_{x=0}^{X-1} 2^x = 3 \times 3 \times 255 = 2295 \text{ samples}
\end{equation}

At stride 8 and 8~kHz, this spans approximately 2.3~seconds of audio.

\subsubsection{Mask Generation}

Skip connections from all $R \times X = 24$ blocks are summed and passed through a $1 \times 1$ convolution producing $C \times N$ channels. After reshaping and applying ReLU activation, the masks are multiplied element-wise with the encoder output.

\subsubsection{Hyperparameter Tuning}

Multiple configurations were explored:

\begin{table}[t]
\centering
\caption{Conv-TasNet Hyperparameter Exploration}
\label{tab:convtasnet_hp}
\begin{tabular}{@{}lccccr@{}}
\toprule
\textbf{Version} & $X$ & $R$ & $B$ & $H$ & \textbf{Val SI-SDR} \\
\midrule
v1 (baseline) & 8 & 3 & 128 & 512 & $-10.64$ dB \\
v2 (scaled) & 10 & 4 & 256 & 1024 & $0.81$ dB \\
v3 (improved) & 8 & 3 & 128 & 512 & $>0$ dB \\
\bottomrule
\end{tabular}
\end{table}

The v2 configuration over-parameterized the model, leading to slow convergence despite improved capacity. Version v3 retained the v1 architecture but increased batch size to 16 and extended training to 200 epochs with learning rate warmup.

% --- Conv-TasNet C++ ---
\subsection{Architecture III: Conv-TasNet (C++ / LibTorch)}
\label{sec:convtasnet_cpp}

To address production deployment requirements, specifically low-latency inference without Python overhead, a full reimplementation of Conv-TasNet was developed in C++ using the LibTorch (PyTorch C++ API) framework.

\subsubsection{Architecture}

The C++ implementation faithfully replicates the Conv-TasNet architecture with parameters $N{=}256$, $L{=}16$, $B{=}128$, $H{=}256$, $P{=}3$, $X{=}8$, $R{=}3$, $C{=}2$, totaling approximately 5.1 million parameters.

\subsubsection{Implementation Details}

Key components implemented in C++:

\textbf{Layer Normalization:} Two variants are provided: \texttt{GlobalLayerNorm} normalizing across all non-batch dimensions, and \texttt{ChannelLayerNorm} normalizing per-channel. Both include learnable affine parameters $\gamma$ and $\beta$.

\textbf{Depth-wise Separable Convolutions:} Implemented using PyTorch's \texttt{groups} parameter set equal to the number of input channels, ensuring per-channel filtering with minimal cross-channel interaction.

\textbf{Loss Functions:} Both time-domain SI-SNR PIT loss and spectral-domain $L_1$ PIT loss are implemented:
\begin{equation}
\mathcal{L}_{\text{spectral}} = \min_{\pi} \sum_{i=1}^{C} \| |\hat{S}_i| - |S_{\pi(i)}| \|_1
\end{equation}

\subsubsection{Audio Processing Pipeline}

The C++ pipeline includes:
\begin{itemize}[leftmargin=*,nosep]
    \item WAV file I/O via \texttt{libsndfile}
    \item Automatic mono conversion (multi-channel averaging)
    \item Resampling to 8~kHz target sample rate
    \item Peak normalization to 0.9 amplitude
    \item Voice Activity Detection (VAD) based silence trimming
    \item STFT computation: $n_{\text{fft}} = 256$, hop length $= 64$, Hann window
\end{itemize}

\subsubsection{Build System}

The project uses CMake with automatic LibTorch and \texttt{libsndfile} detection. Compilation produces a standalone inference binary \texttt{build/inference} that accepts command-line arguments for checkpoint path, input audio file, and output directory.

% --- RCNN ---
\subsection{Architecture IV: Recurrent Convolutional Neural Network (RCNN)}
\label{sec:rcnn}

Unlike the time-domain approaches above, the RCNN operates on STFT magnitude spectrograms, predicting time-frequency masks in the spectral domain.

\subsubsection{STFT Preprocessing}

Input audio is transformed via STFT with $n_{\text{fft}} = 512$, hop length $= 128$, and Hann window, producing a complex spectrogram $\mathbf{X} \in \mathbb{C}^{F \times T}$ where $F = 257$ frequency bins. The model processes only the magnitude $|\mathbf{X}|$, while the mixture phase $\angle \mathbf{X}$ is preserved for reconstruction.

\subsubsection{Encoder}

Three convolutional blocks progressively downsample the frequency dimension:
\begin{align}
\text{Block 1:} &\quad \text{Conv2D}(1 \rightarrow 32, 3{\times}3, \text{stride}(2,1)) \notag \\
\text{Block 2:} &\quad \text{Conv2D}(32 \rightarrow 64, 3{\times}3, \text{stride}(2,1)) \notag \\
\text{Block 3:} &\quad \text{Conv2D}(64 \rightarrow 128, 3{\times}3, \text{stride}(2,1)) \notag
\end{align}

Each block applies batch normalization and ReLU activation. The frequency dimension is reduced from 257 to approximately 32 bins while preserving temporal resolution.

\subsubsection{BiLSTM Core}

The encoder output is reshaped from $(B, 128, F', T)$ to $(B, T, 128 \cdot F')$ and processed by a 2-layer BiLSTM with 256 hidden units per direction (512 total output per time step). Dropout of 0.3 is applied between LSTM layers. 

A linear projection maps the BiLSTM output back to $(128 \cdot F')$ dimensions, and a skip connection adds the encoder output to the projected BiLSTM output before passing to the decoder.

\subsubsection{Decoder}

Three transposed convolutional blocks symmetrically upsample the frequency dimension:
\begin{align}
\text{Block 1:} &\quad \text{ConvTranspose2D}(128 \rightarrow 64, \text{stride}(2,1)) \notag \\
\text{Block 2:} &\quad \text{ConvTranspose2D}(64 \rightarrow 32, \text{stride}(2,1)) \notag \\
\text{Block 3:} &\quad \text{ConvTranspose2D}(32 \rightarrow 1, \text{stride}(2,1)) \notag
\end{align}

Output padding of $(1, 0)$ ensures the reconstructed frequency dimension matches the input (257 bins).

\subsubsection{Mask Application and Waveform Reconstruction}

A final $1 \times 1$ convolution expands the single-channel decoder output to $C$ channels, followed by sigmoid activation to produce soft masks $\mathbf{M}_c \in [0, 1]^{F \times T}$. Separated spectrograms are computed as:
\begin{equation}
\hat{S}_c(f, t) = \mathbf{M}_c(f, t) \cdot |\mathbf{X}(f, t)| \cdot e^{j \angle \mathbf{X}(f, t)}
\end{equation}

The inverse STFT reconstructs the time-domain signals $\hat{s}_c(t)$.

\subsubsection{Inference with Overlapping Segments}

For long audio files, the RCNN employs overlapping segmented inference:
\begin{enumerate}[leftmargin=*,nosep]
    \item Divide the input into segments of 32,000 samples (4~s at 8~kHz)
    \item Apply 50\% overlap between adjacent segments (hop $= 16,000$)
    \item Process each segment independently through the model
    \item Reconstruct via weighted overlap-add, averaging in overlapping regions
\end{enumerate}

This prevents boundary artifacts and enables processing of arbitrarily long recordings.

% ============================================================
\section{Training Pipeline}
\label{sec:training}

\subsection{Distributed Training}

Models are trained using PyTorch's Distributed Data Parallel (DDP) framework across multiple GPUs. The training workflow is:

\begin{enumerate}[leftmargin=*,nosep]
    \item Each GPU receives a unique rank $r \in \{0, \ldots, W{-}1\}$ where $W$ is the world size.
    \item The model is replicated across all GPUs; gradients are averaged via all-reduce after each backward pass.
    \item A \texttt{DistributedSampler} partitions training data to ensure no overlap between GPU data subsets.
    \item Only rank 0 performs logging, checkpoint saving, and metric aggregation.
\end{enumerate}

\subsection{Optimization}

Table~\ref{tab:training_config} summarizes the training configuration shared across models.

\begin{table}[t]
\centering
\caption{Training Hyperparameters}
\label{tab:training_config}
\begin{tabular}{@{}lcccc@{}}
\toprule
\textbf{Parameter} & \textbf{DPRNN} & \textbf{Conv-TN} & \textbf{RCNN} & \textbf{C++} \\
\midrule
Optimizer & Adam & Adam & Adam & Adam \\
Learning Rate & $10^{-3}$ & $10^{-3}$ & $10^{-3}$ & $10^{-3}$ \\
Weight Decay & $10^{-5}$ & 0 & $10^{-5}$ & 0 \\
Batch Size & 8 & 4/GPU & 8 & 4 \\
Max Epochs & 100--500 & 100--200 & 50 & 100 \\
Grad Clip & 5.0 & N/A & 5.0 & 5.0 \\
LR Scheduler & Plateau & Plateau & Plateau & Plateau \\
Sched. Factor & 0.5 & 0.5 & 0.5 & 0.5 \\
Sched. Patience & 5 & 3 & 5 & 3 \\
Mixed Precision & FP16 & BF16 & FP16 & N/A \\
\bottomrule
\end{tabular}
\end{table}

\subsection{Automatic Mixed Precision (AMP)}

Training employs PyTorch's native AMP with \texttt{GradScaler} to accelerate computation. Forward passes execute in FP16 (or BF16 on supported hardware), while gradient accumulation and parameter updates remain in FP32 for numerical stability. This typically provides a 1.5x to 2x speedup with negligible quality impact.

\subsection{Checkpoint Management}

Models save checkpoints containing the model state dict, optimizer state, scheduler state, current epoch, and best validation loss. Due to GitHub's 100~MB file size limit, large checkpoints are automatically split into \texttt{.part0}, \texttt{.part1}, etc., and reassembled at load time:

Each model's API module automatically detects and merges split checkpoints before loading, ensuring seamless deployment without manual intervention.

\subsection{Early Stopping}

Training monitors validation loss (negative mean SI-SNR) and terminates if no improvement is observed for a patience of 5 to 10 epochs. The best-performing checkpoint (by validation loss) is retained for inference.

% ============================================================
\section{Results and Analysis}
\label{sec:results}

\subsection{Evaluation Metrics}

We evaluate separated signals using the following standard metrics:

\begin{itemize}[leftmargin=*,nosep]
    \item \textbf{SI-SNR (dB):} Scale-invariant SNR as defined in Eq.~\eqref{eq:sisnr}. Higher is better; values above 10~dB indicate good separation.
    \item \textbf{SI-SNRi (dB):} SI-SNR improvement over the unprocessed mixture: $\text{SI-SNRi} = \text{SI-SNR}_{\text{est}} - \text{SI-SNR}_{\text{mix}}$.
    \item \textbf{SDRi (dB):} Source-to-Distortion Ratio improvement, computed via \texttt{mir\_eval}. Accounts for spatial distortion, interference, and artifacts.
    \item \textbf{PESQ:} Perceptual Evaluation of Speech Quality, a MOS-correlated metric on a scale of 1.0 (bad) to 4.5 (excellent).
    \item \textbf{STOI:} Short-Time Objective Intelligibility, ranging from 0 to 1, measuring speech intelligibility.
\end{itemize}

\subsection{Quantitative Results}

Table~\ref{tab:results} presents the separation performance across models and speaker configurations.

\begin{table}[t]
\centering
\caption{Speech Separation Performance (SI-SNR in dB)}
\label{tab:results}
\begin{tabular}{@{}lccccc@{}}
\toprule
\textbf{Model} & \textbf{Domain} & \textbf{2-spk} & \textbf{3-spk} & \textbf{5-spk} \\
\midrule
DPRNN-TasNet & Time & 12.50 & \textbf{15.28} & 6.38 \\
Conv-TasNet (Py) & Time & 0.81 & N/A & N/A \\
Conv-TasNet (C++) & Time & \multicolumn{3}{c}{\textit{(production deployment only)}} \\
RCNN & Spectral & 8 to 10 & N/A & N/A \\
\midrule
Mixture (baseline) & N/A & 0.0 & 0.0 & 0.0 \\
\bottomrule
\end{tabular}
\end{table}

\subsection{Analysis of DPRNN-TasNet Results}

The DPRNN-TasNet achieves the strongest results across all configurations. Several key observations emerge:

\textbf{Curriculum Learning Effectiveness:} The 3-speaker model (Phase~2) achieves the highest SI-SNR of 15.28~dB, benefiting from the 2-speaker pre-training that provides a strong initialization for the intra and inter-chunk RNNs. However, continuing the curriculum to 5 speakers degrades performance to 6.38~dB, suggesting that the model's capacity becomes saturated when simultaneously modeling five overlapping speakers.

\textbf{Dual-Path Advantage:} The alternation between intra-chunk (local pattern) and inter-chunk (global context) processing enables DPRNN to capture both fine-grained spectral details and long-range temporal dependencies. The chunk size of $K = 100$ with 50\% overlap provides a favorable balance between local resolution and computational efficiency.

\textbf{Training Dynamics:} With ReduceLROnPlateau scheduling, the learning rate typically decays 3 to 4 times over 100 epochs, with the most significant improvements occurring in the first 30 epochs.

\subsection{Conv-TasNet Analysis}

The Conv-TasNet results highlight the sensitivity of TCN-based architectures to hyperparameter selection:

\textbf{v1 Baseline Instability:} The initial configuration achieved a negative SI-SDR of $-10.64$~dB, indicating that the model was producing outputs worse than the mixture, likely due to insufficient training duration or suboptimal learning rate scheduling.

\textbf{Scaling Challenges (v2):} Increasing the number of blocks ($X = 10$), repeats ($R = 4$), and channel dimensions $(B = 256, H = 1024)$ improved results to 0.81~dB but remained far below the DPRNN baseline. The over-parameterized model required significantly more training compute to converge.

\textbf{Architecture Efficiency:} Despite fewer parameters than DPRNN, Conv-TasNet's fully convolutional separator avoids the sequential bottleneck of RNNs, enabling potentially faster inference. The receptive field of 2,295 samples spans approximately 2.3~seconds, sufficient for capturing most speech patterns.

\subsection{RCNN Spectral-Domain Analysis}

The RCNN's spectral-domain approach offers complementary advantages and limitations:

\textbf{Spectrogram Masking Interpretability:} Operating on magnitude spectrograms enables direct visualization of learned masks, providing insight into which time-frequency regions the model assigns to each speaker.

\textbf{Phase Limitation:} By processing only magnitudes and reusing the mixture phase for reconstruction, the RCNN inherently limits separation quality. In high-overlap regions where both speakers contribute significantly, the mixture phase may not accurately represent either source's true phase, introducing reconstruction artifacts.

\textbf{BiLSTM Bottleneck:} The encoder-BiLSTM-decoder architecture captures temporal context within the spectrogram, but the BiLSTM's sequential nature limits parallelization during training and inference.

\subsection{Performance Degradation with Speaker Count}

Fig.~\ref{fig:degradation} illustrates the consistent performance degradation as the number of speakers increases. This trend is expected: with $C$ speakers, each source occupies on average $1/C$ of the mixture energy, and time-frequency overlap increases combinatorially.

\begin{figure}[t]
\centering
\fbox{\parbox{0.92\columnwidth}{\centering
\small
\textbf{SI-SNR vs. Number of Speakers (DPRNN-TasNet)} \\[4pt]
\begin{tabular}{c|c}
Speakers & SI-SNR (dB) \\ \hline
2 & 12.50 \\
3 & 15.28 \\
5 & 6.38
\end{tabular}
\\[4pt]
\textit{Note: 3-speaker achieves highest SI-SNR due to} \\
\textit{extended training (500 epochs) and curriculum transfer.}
}}
\caption{SI-SNR performance as a function of speaker count for the DPRNN-TasNet model.}
\label{fig:degradation}
\end{figure}

\subsection{Domain Mismatch Analysis}

A critical finding is the significant performance gap between synthetic evaluation mixtures and real-world recordings. Models trained on clean LibriSpeech mixtures exhibit degradation when tested on:

\begin{itemize}[leftmargin=*,nosep]
    \item \textbf{Reverberant audio:} Room acoustics introduce convolutive mixing not present in training data.
    \item \textbf{Non-speech interference:} Background music, mechanical noise, and environmental sounds confuse separation models trained exclusively on speech mixtures.
    \item \textbf{Speaker variability:} Real recordings feature diverse accents, emotions, and speaking styles beyond the read-speech distribution of LibriSpeech.
    \item \textbf{Recording quality:} Microphone characteristics, compression artifacts, and dynamic range differences introduce distribution shift.
\end{itemize}

The MUSAN noise augmentation used in the Conv-TasNet pipeline partially mitigates this issue by exposing the model to diverse noise conditions during training, but reverberant mixing remains unaddressed.

% ============================================================
\section{Comparative Analysis}
\label{sec:comparison}

Table~\ref{tab:comparison} provides a comprehensive comparison across the four architectures along multiple dimensions.

\begin{table*}[t]
\centering
\caption{Comprehensive Comparison of Speech Separation Architectures}
\label{tab:comparison}
\begin{tabular}{@{}lcccc@{}}
\toprule
\textbf{Criterion} & \textbf{DPRNN-TasNet} & \textbf{Conv-TasNet (Py)} & \textbf{Conv-TasNet (C++)} & \textbf{RCNN} \\
\midrule
Operating Domain & Time-domain & Time-domain & Time-domain & Spectral (STFT) \\
Separator Type & Dual-Path BiLSTM & TCN (dilated conv) & TCN (dilated conv) & Encoder-BiLSTM-Decoder \\
Encoder Params ($N$, $L$) & (64, 2) & (512, 16) & (256, 16) & STFT ($n_{\text{fft}}{=}512$) \\
Key Strength & Long-range + local context & Fully parallel separator & Low-latency C++ inference & Interpretable spectral masks \\
Key Weakness & Sequential RNN bottleneck & Hyperparameter sensitivity & Requires C++ compilation & Phase approximation \\
Multi-speaker (2 to 5) & \checkmark & 2 to 3 speakers & 2 speakers & 2 speakers \\
Noise Robustness & No augmentation & MUSAN augmentation & No augmentation & No augmentation \\
Best SI-SNR (2-spk) & 12.50 dB & 0.81 dB & N/A & 8 to 10 dB \\
Training Stability & Stable & Unstable initially & Stable & Stable \\
Long Audio Handling & Chunked overlap-add & Segment-based & Segment-based & Overlapping segments \\
\bottomrule
\end{tabular}
\end{table*}

\subsection{Time-Domain vs.\ Spectral-Domain}

The time-domain models (DPRNN-TasNet, Conv-TasNet) operate directly on raw waveforms through learned encoder-decoder pairs, jointly optimizing the representation and separation. The spectral-domain RCNN operates on hand-crafted STFT features, offering interpretability at the cost of phase approximation errors. Our results suggest that time-domain approaches generally achieve superior SI-SNR when sufficient training is provided, consistent with findings in the literature.

\subsection{Recurrent vs.\ Convolutional Separators}

The DPRNN's bidirectional LSTM separator captures both forward and backward temporal context but processes sequences sequentially. Conv-TasNet's TCN separator enables fully parallel computation during inference, potentially offering faster wall-clock times despite similar parameter counts. The DPRNN's dual-path strategy provides an effective compromise, limiting the sequential processing to chunk length $K$ rather than the full sequence length $T$.

\subsection{Python vs.\ C++ Deployment}

The C++ Conv-TasNet implementation demonstrates the feasibility of production deployment without Python dependencies. Key advantages include:
\begin{itemize}[leftmargin=*,nosep]
    \item Elimination of Python interpreter overhead
    \item Direct memory management via C++ RAII semantics
    \item Standalone binary deployment without environment management
    \item Integration with C/C++ audio processing pipelines
\end{itemize}

The trade-off is increased development complexity and the requirement for compiled LibTorch libraries on the target platform.

% ============================================================
\section{System Architecture}
\label{sec:system}

\subsection{Web Application}

The complete system is deployed as a full-stack web application with the following components:

\textbf{Frontend:} A single-page HTML/CSS/JavaScript application featuring drag-and-drop audio upload, model selection (DPRNN, Conv-TasNet, Conv-TasNet-C++, RCNN), configurable speaker count (2--5), real-time audio preview, spectrogram visualization of separated sources, and downloadable output files. The interface supports dark/light theme switching for user accessibility.

\textbf{Backend:} A FastAPI server exposes a unified \texttt{/api/separate} endpoint that dynamically dispatches to the selected model's API module. The modular architecture uses Python's \texttt{importlib} for dynamic model loading, ensuring that a failure in one model does not prevent the server from serving requests for other models.

\textbf{Model Serving:} Each model implementation exposes a standardized \texttt{separate()} function accepting an audio file, model identifier, and speaker count. The C++ model is invoked via subprocess, executing the compiled inference binary with appropriate command-line arguments.

% ============================================================
\section{Conclusion}
\label{sec:conclusion}

This paper presented a comprehensive comparative study of four deep learning architectures for cocktail party speech separation. The DPRNN-TasNet achieved the best separation quality with an SI-SNR of 15.28~dB for 3-speaker mixtures, demonstrating the effectiveness of dual-path recurrent processing for capturing both local and global temporal structure. Conv-TasNet offered architectural efficiency through fully convolutional separators but required careful hyperparameter tuning to achieve competitive results. The RCNN's spectral-domain approach provided interpretable masks at the cost of phase approximation errors. The C++ Conv-TasNet implementation demonstrated the feasibility of production-grade deployment.

Key findings include: (1)~curriculum learning enables effective multi-speaker generalization but performance degrades beyond 3 to 4 speakers, (2)~time-domain methods outperform spectral methods when sufficient training is provided, (3)~significant domain mismatch exists between synthetic training mixtures and real-world recordings, and (4)~noise augmentation (MUSAN) partially mitigates distribution shift.



% ============================================================
\section*{Acknowledgments}

This work was conducted as part of the Pattern Recognition and Machine Learning (PRML) course project. Computational resources were provided by institutional GPU clusters.



\end{document}
